\section{Related work and contribution}
\label{sec:related}

We reviewed twenty studies from 2023--2026 found by keyword searches on wind downscaling, station exposure and roughness, forest flow, on-road wind measurement and the effect of wind on vehicle energy, together with the classical references they build on. Table~\ref{tab:literature} summarises the most relevant ones by what they predict, what they use as truth, how they validate and whether they reach the vehicle.

\begin{table}[!htbp]
\centering
\caption{Representative studies on the wind near roads and vehicles. Height: height of the predicted wind. Truth: observations against which the method is evaluated. Held out: whether evaluation sites (S) or periods (T) were withheld, or another region (R); --: not done or not reported.}
\label{tab:literature}
\footnotesize
\setlength{\tabcolsep}{3pt}
\begin{tabular}{>{\raggedright\arraybackslash}p{0.16\textwidth}>{\raggedright\arraybackslash}p{0.14\textwidth}>{\raggedright\arraybackslash}p{0.2\textwidth}>{\raggedright\arraybackslash}p{0.2\textwidth}>{\centering\arraybackslash}p{0.07\textwidth}>{\raggedright\arraybackslash}p{0.11\textwidth}}
\toprule
Study & Height, setting & Truth & Method & Held out & Vehicle link \\
\midrule
\citet{DujardinLehning2022} & 10\,m, Alps & meteorological stations & CNN, NWP + terrain & S & -- \\
\citet{Dupuy2023} & 10\,m, SE France & 1\,km WRF fields & CNN, 9\,km $\to$ 1\,km & -- & -- \\
\citet{LeToumelin2024} & 10\,m, Alps & 273 stations & NN correction + CNN downscaling & S, T, R & -- \\
\citet{Lian2024} & near-surface, E China & stations & deep learning, terrain + station links & -- & -- \\
\citet{TroenPetersen1989} & 10--100\,m, wind resource & masts & roughness, obstacle and orography models (WAsP) & S & -- \\
\citet{Herb2007} & along railway tracks & stations & standards-based roughness and orography transfer & -- & rail crosswind risk \\
\citet{Wang2024z0} & roughness length & 101 towers & $z_0$ from tower profiles & -- & -- \\
\citet{Chaudhari2026}; \citet{Topic2026} & canopy to hub height & tower data & LES with forest from laser scanning & -- & -- \\
\citet{McTavish2024} & 1--3\,m on road, Canada & vehicle probes & measurement, video classification & -- & chan\-nel\-ling \\
\citet{Reymert2024,Reymert2025} & vehicle, Norwegian bridges & vehicle anemometer, GNSS/IMU & Kalman filter for crosswind loads & -- & cross\-wind load \\
\citet{Filla2025data,Filla2026} & 6\,m station = vehicle & none for wind & nearest road weather station & -- & drag, 5 years \\
\citet{Askerdal2024crosswind,Askerdal2025winter} & vehicle & CFD; truck anemometer & simplified $C_dA(\beta)$; motion resistance & -- & drag \\
\textbf{This study} & roadside (10\,m) and vehicle (3.4\,m) & \noTrv{} roadside stations (\seTrv{} for development); van anemometer & physical ladder + parsimonious exposure model, four baselines & S, T & drag energy \\
\bottomrule
\end{tabular}
\end{table}

Three gaps stand out. First, downscaling studies predict the wind at \SI{10}{m} over terrain and are evaluated at meteorological stations, which are sited to avoid the sheltered places where roads run; the roadside wind is not their target. Their skill can also be regional: the speed correction of \citet{LeToumelin2024} degraded outside the Alps. Second, roughness and forest studies explain the physics but do not provide a correction that can be evaluated against an independent network. Third, vehicle studies either measure the on-road wind on a few routes \citep{McTavish2024,Reymert2024} or use station wind without correction \citep{Filla2026}, so that the error of the station wind at the road is unknown.

This paper contributes to closing these gaps in four ways. (i) It uses a national road weather network as the truth for the station-to-roadside transfer, with blocked-time $\times$ leave-one-station-out validation over 23 months and a hold-out year, and compares the correction with the alternatives that a practitioner would use: the nearest station rescaled by one constant, the nearest other road weather station, and a 1-km NWP analysis, raw, rescaled by one constant and combined with the exposure model. (ii) It decomposes the transfer into physically interpretable steps using open laser-scanning data and shows which of them matter. (iii) It quantifies the wind penalty in the air-drag energy of a tractor--semitrailer at the station sites, with its uncertainty, and shows how the choice of wind source biases it. (iv) It tests the result at vehicle height and reports where the correction fails, including strong wind and wind direction.
